\documentclass{beamer}
\usepackage[utf8]{inputenc}
\usepackage[portuguese]{babel}
\usepackage{amsfonts,amsmath}
\usepackage{booktabs}
\usepackage{graphicx}
\usetheme{ufms}
\usefonttheme[onlymath]{serif}

\titlebackground*{assets/background}

\title{O ``R'' do RAG}
\subtitle{Recuperação de Artigos sobre IA e Licitações}
\course{Inteligência Artificial — PPGCC / FACOM / UFMS 2026.1}
\author{Luiz Fernando Postingel \textbf{Quirino}}
\IDnumber{Mestrado — Trabalho 1}

\begin{document}
\maketitle
\setcounter{framenumber}{0}

% ============================================================
% SLIDE 1 — O que é + corpus
% ============================================================

\begin{frame}{Sistema de recuperação híbrido}
  \begin{columns}
    \begin{column}{0.55\textwidth}
      \begin{itemize}
        \item Busca artigos sobre \textbf{IA + compras públicas}
        \item Corpus bilíngue construído do zero:
          \begin{itemize}
            \item 3.106 docs (2.135 EN + 971 PT)
            \item OpenAlex + BDTD + arXiv
          \end{itemize}
        \item Embrião do mestrado (módulo de busca)
      \end{itemize}
    \end{column}
    \begin{column}{0.4\textwidth}
      \begin{block}{Achado na coleta}
        arXiv: só 42 on-topic.\\
        Tema \alert{nichado}: muita IA, muita licitação, pouco no cruzamento.
      \end{block}
    \end{column}
  \end{columns}
\end{frame}

% ============================================================
% SLIDE 2 — Recuperadores + módulos
% ============================================================

\begin{frame}{Recuperadores e módulos}
  \textbf{3 recuperadores base:}
  \begin{description}
    \item[BM25] esparso — casa termos exatos (IDF $\approx$ Naïve Bayes)
    \item[KNN/TF-IDF] busca por cosseno no espaço vetorial
    \item[Denso] embeddings multilíngues — casa significado PT$\leftrightarrow$EN
  \end{description}

  \vspace{0.8em}
  \textbf{5 módulos de aprofundamento:}
  \begin{columns}[T]
    \begin{column}{0.5\textwidth}
      \begin{itemize}
        \item M1: re-ranking (Reg.\ Logística)
        \item M2: clustering (K-means)
        \item M3: expansão (regras assoc.)
      \end{itemize}
    \end{column}
    \begin{column}{0.5\textwidth}
      \begin{itemize}
        \item M4: otimização (grid search)
        \item M5: fusão RRF (BM25+denso)
      \end{itemize}
    \end{column}
  \end{columns}
\end{frame}

% ============================================================
% SLIDE 3 — Resultados
% ============================================================

\begin{frame}{Resultados}
  \begin{table}
    \centering\small
    \begin{tabular}{@{}lcccc@{}}
      \toprule
      \textbf{Sistema} & \textbf{P@10} & \textbf{R@10} & \textbf{MAP} & \textbf{nDCG@10} \\
      \midrule
      BM25              & 0,727 & 0,644 & 0,830 & 0,791 \\
      KNN/TF-IDF        & 0,720 & 0,639 & 0,740 & 0,697 \\
      Denso             & 0,240 & 0,172 & 0,180 & 0,306 \\
      RRF (BM25+denso)  & 0,580 & 0,456 & 0,552 & 0,642 \\
      \textbf{M1 (re-ranking)} & \textbf{0,740} & \textbf{0,651} & \textbf{0,864} & \textbf{0,794} \\
      M3 (expansão)     & 0,567 & 0,515 & 0,608 & 0,643 \\
      \bottomrule
    \end{tabular}
  \end{table}

  \vspace{0.3em}
  {\small 15 consultas, 443 julgamentos de relevância (0/1/2).\\
  \textbf{M1 lidera} (MAP 0,864). BM25 é baseline forte (0,830).}
\end{frame}

% ============================================================
% SLIDE 4 — Viés de pooling (o achado honesto)
% ============================================================

\begin{frame}{O achado honesto: viés de pooling}
  \begin{alertblock}{Por que o denso pontua baixo?}
    O gabarito foi construído por pooling lexical (BM25/KNN).\\
    O denso acha relevantes \textbf{fora do pool} $\rightarrow$ marcados como 0.\\
    Não é falha do modelo. É \textbf{viés de pooling} (Zobel, 1998).
  \end{alertblock}

  \vspace{0.5em}
  \textbf{Implicação para o mestrado:}
  \begin{itemize}
    \item Anotar relevância à mão não escala
    \item Motiva aprendizado \textbf{semissupervisionado}
  \end{itemize}
\end{frame}

% ============================================================
% SLIDE 5 — Lições + demo
% ============================================================

\begin{frame}{Lições e acesso ao vivo}
  \textbf{Três lições:}
  \begin{enumerate}
    \item Esparso e denso são \textbf{complementares}
    \item Domínio nichado + bilíngue exige múltiplas fontes
    \item Qualidade da avaliação depende da \textbf{completude do gabarito}
  \end{enumerate}

  \vspace{1em}
  \textbf{Demo ao vivo} (links no LINKS.txt):
  \begin{itemize}
    \item v1: terminal web — pipeline completo
    \item v2: interface web com LLM — busca referências por upload
  \end{itemize}
\end{frame}

% ============================================================
% SLIDE 6 — Encerramento
% ============================================================

\appendix
\backmatter
\end{document}
